\section{The Three Comings of the Lord}\label{sec:three-comings}

\textit{Advent} comes from the Latin \textit{adventus}, meaning a coming or
arrival. In the Vulgate New Testament, it renders the Greek
\textit{parousia}, the word for a royal arrival or Christ's return at the
end of time. Where Matthew reads, ``For as lightning cometh out of the east
and appeareth even into the west: so shall also the coming of the Son of
man be'' (Matt 24:27), the Vulgate has \textit{ita erit et adventus Filii
hominis}.

Preaching at Clairvaux in the twelfth century, St.~Bernard opened his
\textit{First Sermon for Advent} (I.1) by observing how quickly a familiar
name loses its sense:

\begin{quote}
To-day, my brethren, we celebrate the beginning of the season of advent.
The name of this, as of all our other solemnities, is familiar enough and
well known to the entire world. But perhaps the meaning of the name is not
so generally understood. For the unhappy children of Adam, taken up as they
are with the pursuit of things transitory and perishable, pay but little
attention to what is salutary and true.
\end{quote}

What escapes distracted men, in Bernard's Latin, is the \textit{ratio
nominis}, the reason behind the name. The Church keeps Advent not merely as
a countdown to a past birthday, but as a watch before Christ's coming in
the flesh at Bethlehem, His hidden coming into the soul by grace, and His
open coming in glory to judge the living and the dead.

\subsection{Rain on the Fleece and Lightning Across Heaven}

Instructing candidates for Baptism in fourth-century Jerusalem on the
Creed's confession that Christ shall come in glory to judge the living and
the dead, St.~Cyril of Jerusalem set the two visible comings side by side
(\textit{Catechesis} XV.1):

\begin{quote}
We preach not one advent only of Christ, but a second also, far more
glorious than the former. For the former gave a view of His patience; but
the latter brings with it the crown of a divine kingdom. For all things,
for the most part, are twofold in our Lord Jesus Christ: a twofold
generation; one, of God, before the ages; and one, of a Virgin, at the
close of the ages: His descents twofold; one, the unobserved, like rain on
a fleece; and a second His open coming, which is to be. In His former
advent, He was wrapped in swaddling clothes in the manger; in His second,
He covereth Himself with light as with a garment. In His first coming, He
endured the Cross, despising shame; in His second, He comes attended by a
host of Angels, receiving glory. We rest not then upon His first advent
only, but look also for His second.
\end{quote}

Cyril's silent rain echoes Psalm~71~[72]:6: ``He shall come down like rain
upon the fleece; and as showers falling gently upon the earth.'' Later in
the same lecture (\textit{Catechesis} XV.10), pairing that verse with
Matthew 24:27, he contrasts the unseen Nativity with the return from heaven:
``not alone as before, but with many, escorted by tens of thousands of
Angels; nor secretly as the dew on the fleece; but shining forth openly as
the lightning.'' Between them, Cyril adds (\textit{Catechesis} XV.4), ``as
to the time let no one be curious,'' and let none ``supinely slumber.''

\subsection{The Hidden Middle Coming}

How does a Christian standing on an ordinary December weekday pass from
Bethlehem to the Last Day? St.~Bernard took up that question in his
\textit{Fifth Sermon for Advent} (V.1). While the first and last comings
are manifest to all, he teaches, a third stands hidden in the middle, seen
within themselves by the elect alone. Bernard's Latin calls the three
\textit{primus} (first), \textit{medius} (middle), and \textit{ultimus}
(last): Christ came \textit{in carne et infirmitate} (in flesh and
weakness), comes now \textit{in spiritu et virtute} (in spirit and power),
and shall come \textit{in gloria et majestate} (in glory and majesty). (In
the English translation below, ``the second'' numbers this hidden middle
coming, not the Last Day.)

\begin{quote}
Therefore, in the first He came in the flesh and in weakness, in the second
He comes in spirit and in virtue, and in the third He shall come in glory
and majesty. For it is through virtue that glory is reached, as the
Psalmist sings, ``The Lord of virtues, He is the King of glory.'' And
similarly, the same Prophet says in another place, ``In the sanctuary I
have come before Thee, to see Thy virtue and Thy glory.'' Thus, this
intermediate advent is the way, so to speak, by which we must travel from
the first to the last. In the first Christ was our Redemption, in the third
He shall appear as our Life, whilst in this second He is our Repose and our
Consolation, so that we may ``sleep among the midst of lots.''
\end{quote}

Bernard names this middle coming \textit{via qu\ae dam}, a kind of road
between Bethlehem and judgment. In his \textit{Third Sermon for Advent}
(III.4), he summarizes the same triad with three prepositions---\textit{ad
homines, in homines, contra homines}, ``His coming to men, His coming into
men, and His coming against men.'' By His birth He came to all; He comes
into those who receive Him according to John 14:23: ``If any one love me,
he will keep my word. And my Father will love him and we will come to him
and will make our abode with him.''

How, then, is a Christian to ``keep'' that word? In the next paragraph
(\textit{Fifth Sermon for Advent} V.2), Bernard steps straight into the
pantry. Leaving God's word merely on the shelf of memory is no better than
locking a loaf in a chest (\textit{in arca}):

\begin{quote}
Keep the word of God in the same way in which the food of thy body can be
best preserved; for it is a living bread and the food of thy soul. Material
bread, whilst it is kept in the cupboard may be stolen by thieves, or
nibbled by mice, or may become stale and unfit for use. But once it has
been assimilated, we have nothing further to fear from such possibilities.
It is thus we should keep the word of God, for ``blessed are they who keep
it'' so. Receive this spiritual bread, therefore, into the stomach of your
minds, and having been there properly digested, let it be transmitted
thence to the affections and the will. Do not forget (like David) to eat
thy bread, lest thy heart (like his) should wither, but nourish thyself
therewith that thy soul may be filled ``as with marrow and fatness.''
\end{quote}

In Bernard's Latin, the word must pass \textit{in affectiones tuas, et in
mores tuos}---into our affections and into our daily conduct. Then, he adds
(\textit{Fifth Sermon} V.3), the keeping turns around: ``My brethren, if
you thus keep the word of God, there can be no doubt that it, in turn,
shall keep you. For the Son will come to you with the Father, the great
Prophet will come to you to rebuild the Jerusalem of your souls and to make
all things new.''

\subsection{Lifting Up Our Heads Before the Judge}

Medieval liturgical writers carried St.~Bernard's distinction straight into
the sanctuary. The twelfth-century Parisian master John Beleth
(\textit{Summa on the Ecclesiastical Offices}, ch.~64) summarized Christ's
threefold coming (\textit{triplex}) as \textit{in carne} (in the flesh),
\textit{in mente} (day by day in human souls through the Holy Spirit), and
\textit{in majestate} (in majesty at the judgment).

Other commentators noted how the liturgy itself pairs the two visible
advents. Sicard of Cremona (\textit{Mitrale} V.1) spoke of a twofold coming
in humility and majesty, while William Durandus (\textit{Rationale}
VI.2.2--3) added Christ's coming to each soul at death even as the seasonal
Office celebrates Bethlehem and the Last Day.

Because Bethlehem and judgment belong to one Lord, the First Sunday of
Advent opens with the shaking heavens of Luke 21:25--33 and Christ's
command: ``But when these things begin to come to pass, look up and lift up
your heads, because your redemption is at hand'' (Luke 21:28). Preaching in
Rome amid pestilence and sudden disaster, Pope St.~Gregory the Great
(\textit{Homily~1 on the Gospels}, sec.~3, read at Matins of the First
Sunday; Section~\ref{sec:week-1}) urged believers not to cower when the
world shakes: ``To lift up our heads then, is to raise up our minds to the
joys of the heavenly country.''

Even on Christmas Eve, hours before Midnight Mass, the Church refuses to
separate the Infant Redeemer from the returning Judge. On the morning of
24~December, the Collect of the Vigil of the Nativity
(Section~\ref{sec:vigil}; Appendix~\ref{app:prayers}) asks God ``that we,
who now with joy receive Thine only begotten Son as our Redeemer, may
behold Him also without fear, when He cometh as our judge.''
